\documentclass[11pt]{article}

\usepackage[margin=1in]{geometry}
\usepackage[T1]{fontenc}
\usepackage{lmodern}
\usepackage{xcolor}
\usepackage{hyperref}
\usepackage{listings}
\usepackage{longtable}
\usepackage{booktabs}
\usepackage{array}
\usepackage{enumitem}

\hypersetup{
  colorlinks=true,
  linkcolor=blue!55!black,
  urlcolor=blue!65!black,
  pdftitle={FASTEIG User's Guide}
}

\lstdefinestyle{terminal}{
  basicstyle=\ttfamily\small,
  backgroundcolor=\color{black!3},
  frame=single,
  framerule=0.3pt,
  columns=fullflexible,
  keepspaces=true,
  breaklines=true,
  showstringspaces=false
}
\lstset{style=terminal}

\newcommand{\fasteig}{\textsc{fasteig}}
\newcommand{\code}[1]{\texttt{#1}}

\title{\fasteig{} User's Guide}
\author{Based on the FASTEIG source distribution}
\date{Repository version 1.0.0\\Guide revised August 7, 2026}

\begin{document}
\setlength{\emergencystretch}{3em}
\maketitle
\tableofcontents
\newpage

\section{Introduction}

\fasteig{} is a C++17 library and example program for computing selected
eigenpairs of large sparse, real matrices.  A stored standard problem has the
form
\[
    A x = \lambda x,
\]
and an optional generalized problem has the form
\[
    A x = \lambda B x.
\]
The command-line driver reads the matrices and solver settings from a small
two-column input file.  It can run the supported algorithms on a CPU; CUDA
builds can additionally execute Lanczos and Block Lanczos on an NVIDIA GPU.

This guide describes the repository as it currently exists.  In particular,
Section~\ref{sec:cmake-limitations} distinguishes the CMake library build from
the Makefile-based example and CUDA builds.

\section{Prerequisites}

\subsection{Required software}

\begin{itemize}[leftmargin=*]
  \item A C++17 compiler.  The top-level CMake file recognizes GCC, Clang,
        Intel LLVM (\code{icpx}), and Intel Classic compilers.
  \item CMake 3.20 or newer for the CMake library build.
  \item Eigen headers.  Eigen 3.4 is a suitable choice.  Eigen is header-only,
        but its include directory must contain \code{Eigen/Core}.
  \item A BLAS/LAPACK implementation.  The CMake build can use system
        BLAS/LAPACK, Intel oneMKL, Cray LibSci, or a fetched OpenBLAS copy.
\end{itemize}

\subsection{Optional CPU packages}

\paragraph{Intel oneMKL.}
Enable oneMKL with \code{-DUSE\_MKL=ON}.  FASTEIG deliberately uses the MKL
LP64 integer interface and the GNU threading layer.  Set \code{MKLROOT} either
in the environment (usually by sourcing the oneAPI setup script) or on the
CMake command line.  MKL also supplies Pardiso, which FASTEIG uses for CPU
shift-and-invert (\code{which SAI}) when Eigen's MKL integration is enabled.

\paragraph{Cray LibSci.}
On a Cray system, configure with \code{-DUSE\_LIBSCI=ON} and use the Cray
compiler wrappers.  The CMake logic looks first for \code{sci\_gnu}, then for
the \code{cray-libsci} pkg-config package.  \code{USE\_MKL} and
\code{USE\_LIBSCI} are mutually exclusive.

\paragraph{OpenBLAS.}
With neither MKL nor LibSci selected, CMake searches for system BLAS and
LAPACK.  If they are unavailable, it fetches OpenBLAS 0.3.27 into
\code{openblas/} in the source tree and builds it.  This fallback requires Git
and network access during configuration.  Installing BLAS/LAPACK through the
operating system avoids that download.

\subsection{Optional GPU packages}

A GPU build requires:
\begin{itemize}[leftmargin=*]
  \item an NVIDIA GPU supported by the selected CUDA toolkit;
  \item the CUDA compiler \code{nvcc}; and
  \item cuSPARSE, cuSOLVER, cuBLAS, and the CUDA runtime.
\end{itemize}

NVIDIA cuDSS is optional for ordinary GPU Lanczos calculations, but is
required for GPU shift-and-invert calculations (\code{which SAI}).  The
repository supplies \code{tools/install\_cudss.sh}, which can install a CUDA
12 or CUDA 13 Python wheel into project space, or unpack an NVIDIA tarball.

\section{Building with CMake}
\label{sec:cmake}

Run all commands in this section from the repository root.

\subsection{Basic CPU build}

If Eigen is already installed, point \code{EIGEN\_DIR} at the directory that
directly contains \code{Eigen/Core}.  For example:

\begin{lstlisting}
cmake -S . -B build \
  -DCMAKE_BUILD_TYPE=Release \
  -DEIGEN_DIR=/path/to/eigen-3.4.0 \
  -DBUILD_EXAMPLES=OFF \
  -DBUILD_TESTS=OFF
cmake --build build -j 8
\end{lstlisting}

If \code{EIGEN\_DIR} is omitted, CMake first checks for an existing
\code{eigen/Eigen/Core} below the source directory.  Otherwise it fetches
Eigen's \code{master} branch into \code{eigen/}; this also requires Git and
network access.

The resulting static library is
\code{build/lib/libfasteig\_Generic.a}.  Change its platform suffix with, for
example, \code{-DPLAT=Perlmutter}, which produces
\code{libfasteig\_Perlmutter.a}.

\subsection{CMake build with Intel oneMKL}

After making the oneAPI environment available, use:

\begin{lstlisting}
source /path/to/oneapi/setvars.sh
cmake -S . -B build-mkl \
  -DCMAKE_BUILD_TYPE=Release \
  -DEIGEN_DIR=/path/to/eigen-3.4.0 \
  -DUSE_MKL=ON \
  -DMKLROOT="$MKLROOT" \
  -DBUILD_EXAMPLES=OFF \
  -DBUILD_TESTS=OFF
cmake --build build-mkl -j 8
\end{lstlisting}

If the environment has not set \code{MKLROOT}, replace the shell expansion in
the command with the actual MKL installation prefix.  Do not select the ILP64
interface: the current Eigen/Pardiso path expects LP64 integers.

\subsection{CMake build with Cray LibSci}

A typical Cray configuration is:

\begin{lstlisting}
CC=cc CXX=CC cmake -S . -B build-libsci \
  -DCMAKE_BUILD_TYPE=Release \
  -DEIGEN_DIR=/path/to/eigen-3.4.0 \
  -DUSE_LIBSCI=ON \
  -DPLAT=Perlmutter \
  -DBUILD_EXAMPLES=OFF \
  -DBUILD_TESTS=OFF
cmake --build build-libsci -j 8
\end{lstlisting}

The exact environment modules and compiler wrappers are site-specific.

\subsection{Current CMake limitations}
\label{sec:cmake-limitations}

The current top-level CMake configuration builds the CPU static library only.
It does not yet define CUDA sources or a CUDA option.  It also does not define
an \code{install()} rule.  Therefore ``installation'' currently means retaining
the library in the build tree, or manually copying the archive and the
\code{include/} tree to an application prefix.

Although \code{BUILD\_EXAMPLES} defaults to \code{ON}, this repository revision
does not contain \code{examples/CMakeLists.txt}; similarly, there is no
\code{tests/CMakeLists.txt}.  Use \code{BUILD\_EXAMPLES=OFF} and
\code{BUILD\_TESTS=OFF} for a working CMake configuration.  Build the supplied
command-line programs with the Makefiles as described next.

\section{Building the Command-Line Programs}

The checked-in Makefiles produce the complete runnable examples and are the
current build route for CUDA.  They use settings from \code{make.inc}.  That
file contains machine-specific paths, so make a local copy and edit it rather
than relying on its checked-in values:

\begin{lstlisting}
cp make.inc make.inc.local
# Edit make.inc.local: set CXX, CC, PLAT, EIGEN, MKLROOT,
# CUDA_PATH, CUDA_ARCH, and related flags for this machine.
\end{lstlisting}

Here \code{EIGEN} is the directory containing \code{Eigen/Core},
\code{CUDA\_ARCH} is the NVIDIA compute capability without a decimal point
(for example, \code{80} for SM 8.0), and \code{PLAT} determines the library
file suffix.

\subsection{CPU executable}

For the default MKL/Pardiso configuration:

\begin{lstlisting}
make MAKEINC=make.inc.local
make -C examples MAKEINC=../make.inc.local
\end{lstlisting}

This produces the library \code{libfasteig\_<PLAT>.a} at the repository root
and the executable \code{examples/runfasteig\_cpu}.

To use Eigen's native sparse LDLT instead of MKL/Pardiso, remove MKL flags from
the build consistently:

\begin{lstlisting}
make -C src MAKEINC=../make.inc.local clean
make -C examples MAKEINC=../make.inc.local clean
make MAKEINC=make.inc.local USE_MKL=""
make -C examples MAKEINC=../make.inc.local USE_MKL="" MKL=""
\end{lstlisting}

\subsection{GPU executable}

After setting \code{CUDA\_PATH}, \code{CUDA\_ARCH}, and \code{NVCC} in the
local make include, build with:

\begin{lstlisting}
make MAKEINC=make.inc.local USE_CUDA=1
make -C examples MAKEINC=../make.inc.local USE_CUDA=1
\end{lstlisting}

This creates both \code{runfasteig\_cpu} and \code{runfasteig\_gpu}.  The GPU
driver selects a CUDA executor for \code{solver Lanczos},
\code{solver BlockLanczos}, \code{solver KrylovSchur},
\code{solver BlockKrylovSchur}, and \code{solver Davidson}.  Other solver names
in the GPU executable continue to use their CPU implementations, so use
\code{runfasteig\_cpu} for clarity when selecting them.

The CUDA backend requires a matrix stored in one of FASTEIG's input formats.
Matrix-free operators are CPU-only in the current driver.  Generalized
problems should also be run on the CPU; they are currently supported only by
Krylov--Schur, which is not dispatched to the CUDA backend.

\subsection{Installing and enabling cuDSS}
\label{sec:cudss}

For GPU \code{SAI}, install cuDSS into the default checkout-local location:

\begin{lstlisting}
tools/install_cudss.sh
make -C examples MAKEINC=../make.inc.local USE_CUDA=1
\end{lstlisting}

If direct package downloads are unavailable, download the Linux cuDSS archive
from NVIDIA and run:

\begin{lstlisting}
tools/install_cudss.sh --prefix /project/path/cudss \
  --tarball /path/to/cudss-linux-archive.tar.xz
source /project/path/cudss/cudss-env.sh
make -C examples MAKEINC=../make.inc.local USE_CUDA=1 \
  CUDSS_PATH="$CUDSS_PATH" CUDSS_LIBDIR="$CUDSS_LIBDIR"
\end{lstlisting}

The Makefile enables cuDSS automatically when it finds both
\code{include/cudss.h} and \code{libcudss.so} (or
\code{libcudss.so.0}).  To build CUDA support without cuDSS, pass
\code{USE\_CUDSS=0}; ordinary \code{SR}, \code{LR}, \code{SM}, \code{LM}, and
\code{BE} Lanczos runs do not require cuDSS.

\section{Preparing Matrices and an Input File}

\subsection{Matrix requirements and formats}

The command-line driver reads real, square sparse matrices.  Eigensolvers in
this distribution are intended for symmetric/self-adjoint problems.  The
recommended interchange format is Matrix Market coordinate format
(\code{.mtx}).  A minimal symmetric example is:

\begin{lstlisting}
%%MatrixMarket matrix coordinate real symmetric
3 3 5
1 1 2.0
2 1 -1.0
2 2 2.0
3 2 -1.0
3 3 2.0
\end{lstlisting}

Matrix Market indices are one-based.  For a file declared \code{symmetric} or
\code{skew-symmetric}, store one triangle plus the diagonal; the loader keeps
one triangle and applies the symmetry during matrix-vector products.  A real
coordinate file is safest: the current reader expects row, column, and one
floating-point value per entry.

The parser also accepts:
\begin{itemize}[leftmargin=*]
  \item \code{.rsa}: a Harwell--Boeing/RSA-style text matrix in compressed
        sparse column form; and
  \item \code{.fei}: FASTEIG's binary CSC form, containing three
        \code{FASTEIG\_INT} header values (rows, columns, nonzeros), followed by
        the zero-based column-pointer array, row-index array, and
        \code{FASTEIG\_DOUBLE} values.
\end{itemize}

The extension \code{.hdf} is recognized by format detection but is not loaded
by the current input parser.  Do not use it for command-line runs.  For a
generalized problem, \code{amatrix} and \code{bmatrix} must have the same
format and dimension; \(B\) should be symmetric positive definite.

\subsection{Input-deck syntax}

Each non-comment record contains exactly a key and a whitespace-separated
value.  Keys are case-insensitive.  A line whose first token begins with
\code{\#} or \code{\%} is a comment.  Avoid inline comments and paths containing
spaces.  Matrix paths are interpreted relative to the directory from which the
executable is launched, not relative to the input file itself.

A practical standard-problem input file is:

\begin{lstlisting}
# Matrix and output
amatrix ./Sparse_Matrices/SiO.mtx
output_prefix SiO
output\_format text

# Requested eigenpairs
nev 5
which SR
solver Lanczos

# Iteration controls
maxiter 200
maxdim 20
blksize 5
tol 1e-8
seed 1

# Console output
verbosity 2
print\_every 10
\end{lstlisting}

For a generalized problem, add a matching matrix:

\begin{lstlisting}
amatrix /data/A.mtx
bmatrix /data/B.mtx
solver KrylovSchur
which SR
nev 10
maxiter 500
maxdim 40
tol 1e-10
output_prefix generalized
output_format text
\end{lstlisting}

\subsection{Input keys}

\renewcommand{\arraystretch}{1.15}
\begin{longtable}{>{\ttfamily}p{0.18\textwidth}p{0.18\textwidth}p{0.55\textwidth}}
\toprule
\normalfont Key & Default & Meaning \\
\midrule
\endhead
amatrix & required & Path to the standard matrix \(A\). \\
bmatrix & none & Optional matrix \(B\) for \(Ax=\lambda Bx\).  Currently CPU Krylov--Schur only. \\
solver & Lanczos & Lanczos, BlockLanczos, BlockKrylovSchur, KrylovSchur, Davidson, LOBPCG, Chebyshev, SPPC, RMM-DIIS, or auto. \\
nev & 100 & Number of requested eigenpairs; must be positive and no greater than the matrix dimension. \\
which & LM & Target: SR (smallest algebraic), LR (largest algebraic), SM (smallest magnitude), LM (largest magnitude), BE (both ends), or SAI (shift-and-invert).  Not every solver accepts every target. \\
shift & 0.0 & Shift \(\sigma\) for SAI, which applies an operator based on \((A-\sigma I)^{-1}\). \\
maxiter & 10 & Maximum solver iterations or Krylov steps, depending on the algorithm. \\
maxdim & 20 & Maximum subspace dimension for applicable solvers; SPPC uses it as its split dimension. \\
blksize & 5 & Initial block size for block algorithms. \\
tol & $10^{-6}$ & Positive convergence tolerance. \\
seed & unspecified & Nonnegative random seed for reproducible initial vectors. \\
logfile & logfile & Output prefix; \code{output\_prefix} is an alias. \\
output\_format & text & Output format.  Only \code{text} is currently implemented. \\
verbosity & 1 & Console level 1 prints the run header and final eigenpairs; level 2 also enables iteration summaries.  Values are clamped to 1 or 2. \\
print\_every & 0 & Progress-printing cadence for solvers that use it; zero disables cadence-based progress. \\
verbose & --- & Backward-compatible Boolean key that enables normal console output. \\
\bottomrule
\end{longtable}

The data structure contains a Chebyshev \code{degree} setting (default 10) and
a Davidson \code{maxblks} setting, but the current file parser does not assign
those two keys.  Consequently, command-line input decks use their compiled
defaults unless the C++ driver is changed to set them programmatically.

\subsection{Choosing a solver and target}

For a first run, use \code{Lanczos} with \code{SR} or \code{LR}.  The Lanczos
family accepts \code{SR}, \code{LR}, \code{SM}, \code{LM}, and \code{BE}; scalar
Lanczos also accepts \code{SAI}.  \code{KrylovSchur} supports standard and
generalized problems, but generalized \code{SAI} is rejected.  The block
Krylov--Schur implementation does not support generalized problems or
\code{SAI}.  Davidson accepts \code{SR}, \code{LR}, \code{SM}, \code{LM}, and
\code{BE} for standard problems, but not \code{SAI}.  Chebyshev filtering is
designed for an end of the spectrum; the implementation reverses the operator
for \code{LR}/\code{LM}.

Use \code{SAI} to find eigenvalues near \code{shift}.  It incurs a sparse
factorization setup cost and needs a stored matrix.  The CPU backend reports
either \code{MKL PardisoLDLT} or \code{Eigen SimplicialLDLT}; the GPU backend
reports \code{NVIDIA cuDSS}.

\section{Running on CPUs}

The input paths in the supplied \code{examples/SiO.input} assume that the
working directory is \code{examples}.  Run:

\begin{lstlisting}
cd examples
./runfasteig_cpu SiO.input
\end{lstlisting}

The driver exits with status 0 on success, 1 for a command-line or input error,
and 2 when the solver returns a nonzero status.  With
\code{output\_prefix SiO} and \code{solver Lanczos}, the text log is named
\code{SiO\_Lanczos.log}.

Control CPU threading with the OpenMP and math-library environment.  For an
MKL build on a 16-core allocation:

\begin{lstlisting}
export OMP_NUM_THREADS=16
export MKL_NUM_THREADS=16
export OMP_PLACES=cores
export OMP_PROC_BIND=spread
export OMP_DYNAMIC=false
export MKL_DYNAMIC=false
./runfasteig_cpu SiO.input
\end{lstlisting}

On a Slurm system, request one process and assign its cores explicitly, for
example:

\begin{lstlisting}
srun -N 1 -n 1 -c 16 --cpu-bind=cores \
  ./runfasteig_cpu SiO.input
\end{lstlisting}

FASTEIG's example driver is a single-process program; do not launch multiple
MPI ranks unless each rank is intentionally running an independent problem.

\section{Running on NVIDIA GPUs}

Use the same input deck and select the backend with the executable name.  Keep
a CUDA-enabled solver name such as \code{solver Lanczos},
\code{solver BlockLanczos}, or \code{solver Davidson} in the deck:

\begin{lstlisting}
cd examples
./runfasteig_gpu SiO.input
\end{lstlisting}

For an ordinary Lanczos target such as \code{SR}, no cuDSS installation is
needed.  For interior eigenvalues, use both \code{which SAI} and a shift:

\begin{lstlisting}
solver Lanczos
which SAI
shift 0.25
nev 5
\end{lstlisting}

That run requires a build with \code{USE\_CUDA=1 USE\_CUDSS=1}.  If the program
reports that CUDA SAI requires cuDSS, rebuild after following the cuDSS steps
in Section~\ref{sec:cudss}.

A typical Slurm GPU launch is:

\begin{lstlisting}
export OMP_NUM_THREADS=16
export MKL_NUM_THREADS=16
export OMP_PLACES=cores
export OMP_PROC_BIND=close
export OMP_DYNAMIC=false
export MKL_DYNAMIC=false
srun -n 1 -c 16 --gpus-per-task=1 --gpu-bind=single:1 \
  --cpu-bind=cores ./runfasteig_gpu SiO.input
\end{lstlisting}

GPU work is not selected merely by compiling with CUDA: run the
\code{runfasteig\_gpu} executable and choose a CUDA-enabled solver.  The current
CUDA executor accepts stored CSR data derived from the input matrix; it does
not accept the matrix-free callback path.

\section{Verification}

The repository includes a Python 3 integration harness and a reference spectrum
for \code{SiO.mtx}.  After building \code{examples/runfasteig\_cpu}, run from
the repository root:

\begin{lstlisting}
python3 tests/scripts/run_suite.py --matrix SiO
\end{lstlisting}

To test one solver and one end of the spectrum:

\begin{lstlisting}
python3 tests/scripts/run_case.py \
  --solver Lanczos --which SR --matrix SiO
\end{lstlisting}

The test harness checks reported residuals and eigenvalues against
\code{tests/reference/SiO.out.121mpi}.  It requires only Python 3, not SciPy.

\section{Troubleshooting}

\begin{description}[style=nextline]
  \item[CMake reports that \code{Eigen/Core} is missing.]
    Set \code{EIGEN\_DIR} to the directory immediately above \code{Eigen}, not
    to the \code{Eigen} directory itself.  Common choices are
    \code{/usr/include/eigen3} and a source checkout such as
    \code{/opt/eigen-3.4.0}; verify that
    \code{<EIGEN\_DIR>/Eigen/Core} exists.

  \item[CMake tries to clone OpenBLAS.]
    Install development packages for both BLAS and LAPACK, provide a discoverable
    vendor implementation, or configure with MKL/LibSci.  Otherwise network
    access is expected for the fallback.

  \item[CMake fails while adding \code{examples} or \code{tests}.]
    Reconfigure with \code{-DBUILD\_EXAMPLES=OFF -DBUILD\_TESTS=OFF}.  Build the
    drivers with their Makefile.

  \item[The matrix file cannot be opened.]
    Resolve \code{amatrix} and \code{bmatrix} from the launch directory.  For
    example, the supplied relative path works after \code{cd examples}, not
    necessarily from the repository root.

  \item[A generalized problem is rejected.]
    Use \code{solver KrylovSchur} and the CPU executable.  Ensure \(A\) and
    \(B\) have identical dimensions and file-format extensions.

  \item[GPU execution says a stored CSR matrix is required.]
    Supply \code{amatrix} in \code{.mtx}, \code{.rsa}, or \code{.fei} format.
    The current GPU driver cannot use a matrix-free callback.

  \item[GPU SAI says cuDSS is unavailable.]
    Confirm that \code{CUDSS\_PATH/include/cudss.h} and a cuDSS shared library
    under \code{CUDSS\_LIBDIR} exist, then clean and rebuild with
    \code{USE\_CUDA=1 USE\_CUDSS=1}.  Source the generated
    \code{cudss-env.sh} before both building and running when using a custom
    prefix.

  \item[Performance is lower than expected.]
    Use a Release build, match thread counts to the allocated cores, disable
    dynamic OpenMP/MKL resizing, and choose the correct \code{CUDA\_ARCH} for a
    GPU build.  Avoid oversubscribing a batch allocation.
\end{description}

\section{Quick Reference}

\paragraph{CPU.}
\begin{lstlisting}
make MAKEINC=make.inc.local
make -C examples MAKEINC=../make.inc.local
cd examples
./runfasteig_cpu SiO.input
\end{lstlisting}

\paragraph{GPU.}
\begin{lstlisting}
make MAKEINC=make.inc.local USE_CUDA=1
make -C examples MAKEINC=../make.inc.local USE_CUDA=1
cd examples
./runfasteig_gpu SiO.input
\end{lstlisting}

\paragraph{Compile this guide.}
From the repository root:
\begin{lstlisting}
pdflatex -output-directory=docs docs/fasteig_user_guide.tex
pdflatex -output-directory=docs docs/fasteig_user_guide.tex
\end{lstlisting}

\end{document}

